

\section{Proofs for Signature Scheme in \secref{sec:sig_scheme}}
\label{sec:proof:sig_scheme}
\subsection{Proof of Unforgeability}
\label{app:proof_unforge}
 
\begin{theorem}
	\label{thm:unforge}
	The signature scheme \sigScheme is unforgeable under chosen message attacks, that is assuming that 
	strong QR-RSA problem is hard to solve in group \qrGroup. 		

	\[
	\ufcmaAdv \le \frac{\oracleQueries^2 \log \rsaMod}{\rsaMod} + \oracleQueries \rsaAdv 
	\]
	
        Here \oracleQueries is the number of distinct queries to $\signingAlg[\privKey](\cdot)$
        %\mkr{It should be the number of queries to the sign oracle.
         % BTW, has the notation \sigScheme been defined?}
	
\end{theorem}

\begin{proof}
We will show that if there is an adversary \ufcmaAdversary that can forge a signature in time 
\timeBound, then there is an adversary \rsaAdversary that solves the strong QR-RSA problem with 
with probability \rsaAdv in time \timeBound.
The reduction works as follows. Consider that 
\rsaAdversary simulates \ufcmaAdversary's queries to the 
signing oracle, and the random oracle \primeRO. Let $\genericSetVar[1], \dots, 
\genericSetVar[\numberOfSetsQueried]$  be all the sets that are queried by \ufcmaAdversary 
to the signing oracle. Assume that \oracleQueries is the total number of queries to the random 
oracle that includes queries by the signing oracle when responding to \ufcmaAdversary, and 
queries directly by \ufcmaAdversary. Now assume that \ufcmaAdversary forges a signature on 
some value \genericFnOutput that is not in any of the sets $\genericSetVar[1], \dots, 
\genericSetVar[\numberOfSetsQueried]$. This happens under two conditions.


\enumListStart
	\item There is some $\genericFnOutputAlt \in \bigcup\limits_{\setIdx \in \nats[i]} \genericSetVar[i]$ such that $\primeRO(\genericFnOutputAlt) = 
	\primeRO(\genericFnOutput)$.
	
	\item For all  $\genericFnOutputAlt \in \genericSetVar[1], \dots, 
	\genericSetVar[\numberOfSetsQueried]$  $\primeRO(\genericFnOutputAlt) \neq
	\primeRO(\genericFnOutput)$.
\enumListEnd

The first case happens with probability at most $\frac{\oracleQueries^2 \log \rsaMod}{\rsaMod}$ 
because \primeRO maps to prime numbers in $[1, \rsaMod]$ and there are roughly 
 $\frac{\rsaMod}{\log \rsaMod}$ prime numbers less than $\rsaMod$. 
%\mkr{Should the result be stated using $\lesssim$?}

For the second case, consider the following simulation

%  {\color{violet}[Arpad: corrected the value found the the QR-RSA adversary and called the programmed oracle $\primeRO'$ so it's not mistaken for the real world oracle.]}
\enumListStart
%	\item Upon getting the tuple $(\rsaMod, \genericFnOutput)$ where 
%	$\genericFnOutput \in \qrGroup$, \rsaAdversary first checks that $\genericFnOutput$ is 
%	a generator of $\qrGroup$ by verifying that $\gcd(\rsaMod, \genericFnOutput - 1)  = 
%	1$~\cite{micciancio2005rsa}. Otherwise, \rsaAdversary aborts. 
	
	
	\item \rsaAdversary selects \oracleQueries random prime numbers less than $\rsaMod$, 
	$\randPrimes[1], \dots, 
	\randPrimes[\oracleQueries]$ and computes $\randPrimeAccSet \gets 
	\prod\limits_{\setIdx \in 
	[1, \oracleQueries]} \randPrimes[\setIdx]$. Then, it computes 	$\simFnOutput \gets 
	\genericFnOutput^{\randPrimeAccSet} \bmod \rsaMod$. \rsaAdversary publishes 
	$\simFnOutput$ as part of the public key. 
	
	\item\rsaAdversary selects a random index $\randIdx \in [1, \oracleQueries]$ and programs the 
	random oracle $\primeRO$ such for the $\setIdx$-th oracle query for all $\setIdx \neq \randIdx$, the random 
	oracle returns $\randPrimes[\setIdx]$. For the $\randIdx$-th query, the random oracle returns a 
	random prime number $\rsaExponent \notin \randPrimes[1], \dots, 
	\randPrimes[\oracleQueries]$
	
	%, for $\setIdx \neq \randIdx$ and returns \rsaExponent for the $\randIdx$-th  query. 
	
	\item \rsaAdversary answers the signing oracle queries from \ufcmaAdversary for the sets 
	$\genericSetVar[1], \dots, \genericSetVar[\numberOfSetsQueried]$. Specifically, given some set 
	$\genericSetVar$, \rsaAdversary signs \genericSetVar by computing $\randPrimeAccSet' \gets \{ \randPrimes[i]:\randPrimes[i] \in \randPrimes[1], \dots, \randPrimes[\oracleQueries],  \forall \genericSetElmt \in \genericSetVar, \primeRO(\genericSetElmt)   \neq \randPrimes[i] \}$ 
and returning $\dhVar[1] \gets \genericFnOutput^{P'} = 
	\simFnOutput^{ \left(\frac{1}{\challengeSetAcc}\right)}$

%and sends the challenge 
%	$\tupleDefn{\dhVar[1], \genericRandAltAlt}$ to \indsrAdversary{3} where $\genericRandAltAlt \sample 
%		\setDefn{0,1}^{\secParam}$. 




%\rsaAdversary returns $\simFnOutput^{\left(\frac{1}{\challengeSetAcc}\right)} 
%	\bmod \rsaMod$
	
	
	\item  If \ufcmaAdversary forges a signature on $\genericSetVar'$ such that there is some 
	$\genericSetElmt \in \genericSetVar'$ but $\genericSetElmt \notin \genericSetVar[1] \cup \dots  
	\cup \genericSetVar[\numberOfSetsQueried]$ and $\primeRO'(\genericSetElmt) = \rsaExponent$, 
	then  \rsaAdversary learns some value 
	$\genericVal$ such that 
	$\genericVal^{\rsaExponent}  = 	
	\genericFnOutput^{\randPrimeAccSet} \bmod \rsaMod$. Since, $\rsaExponent$ is co-prime to 
	\randPrimeAccSet, we can find the values of $\euclideanConst[1], \euclideanConst[2]$ using the 
	extended Euclid algorithm such that 
	$\euclideanConst[1] \rsaExponent + \euclideanConst[2] \randPrimeAccSet = 1$. Then, we have 
	$(\genericVal^{\euclideanConst[2]} \genericFnOutput^{\euclideanConst[1]})^{\rsaExponent} 
	= \genericFnOutput$. \rsaAdversary outputs $\genericVal^{\euclideanConst[2]} 
	\genericFnOutput^{\euclideanConst[1]}$ as the $\rsaExponent$-th root of $\genericFnOutput$ 
	
%	\item If \ufcmaAdversary forges a signature on $\genericSetVar'$ such that there is some 
%	$\genericSetElmt \in \genericSetVar'$ such that $\primeRO(\genericSetElmt) = \rsaExponent$, 
%	then \rsaAdversary learns some value $\genericFnInput$ such that 
%	$\genericFnInput^{\rsaExponent} = \simFnOutput = 	
%	\genericFnOutput^{\randPrimeAccSet} \bmod \rsaMod$. Since, \rsaExponent 
	
\enumListEnd

To see that this simulation is indistinguishable first note that in the real 
experiment when \ufcmaAdversary queries the set 
\genericSetVar, its view includes $\qrGen^{1/\accSet_{\setIdx}}$ 
where $\accSet[\setIdx] \gets \prod\limits_{\genericSetElmt \in \genericSetVar[\setIdx]} 
\primeRO(\genericSetElmt)$ and \qrGen is a generator of \qrGroup. In the simulation, 
\ufcmaAdversary's view includes $\simFnOutput^{\left(\frac{1}{\challengeSetAcc}\right)}$. 
Since $\genericFnOutput$ is selected randomly from 
\qrGroup, the probability that it is not a generator is negligible in $\secParam$. Consequently, $\simFnOutput 
\gets \genericFnOutput^{\randPrimeAccSet}$ is also a generator with overwhelming 
probability, i.e., as long as \randPrimeAccSet is co-prime to $\sgPrime[1]$ and $\sgPrime[2]$. From 
the above, we have that the signature in the simulation is valid and 
indistinguishable to \ufcmaAdversary. Finally, $\rsaAdversary$ solves the strong QR-RSA problem 
using the simulation when it programs the random oracle such that $\primeRO(\genericSetElmt) = 
\rsaExponent$, this happens with probability at least 
$1/\oracleQueries$. Thus, if $\ufcmaAdversary$ is successful 
in the forging a signature in the simulation with probability \ufcmaAdv, then \rsaAdversary solves the 
strong QR-RSA problem also with probability at least $\rsaAdv/\oracleQueries$. 
\end{proof}


\subsection{Proof of Indistinguishable Signature under Randomization (IND-SR)}
\label{app:proof_indsr}

%{\color{violet}[Arpad: made similar corrections to the ones in the uf proof]}
\begin{theorem}
\label{thm:indsr}
	The signature scheme \sigScheme produces indistinguishable signatures under randomization, that is 
	assuming the strong QR-RSA problem is hard to solve in \qrGroup where $\rsaMod$ is the product of safe primes.
	
	\[
	\Advantage{\indsrLabel}{\sigScheme}(\timeBound, 
	\oracleQueries, \oracleQueries') \le   \frac{\oracleQueries' }{\setSize{\qrGroup}}  + \oracleQueries 
	\rsaAdv
	\]
Here \oracleQueries is the number of distinct queries to the random oracle \primeRO and 
$\oracleQueries'$ is the number of distinct queries to the random oracle \ro.


\end{theorem}


\begin{proof}


Consider that \indsrAdversary{3} gets a challenge $\langle \qrGen^{\left(\frac{\genericRand \primeRO(\genericSetElmt[\indBit])}{\accSet}\right)} \bmod \rsaMod, \ro(\qrGen^{\genericRand} \bmod \rsaMod)\rangle$ where $\qrGen^{\frac{ 
	1}{\accSet}} \bmod \rsaMod$ is a signature on the set $\genericSetVar$. 
Since $\ro$ is a random oracle and $\qrGen^{\genericRand} \bmod \rsaMod$ is uniformly distributed in 
\qrGroup \cite{cramer2000signature}, \indsrAdversary{3} gets $\qrGen^{\genericRand} \bmod 
\rsaMod$ from 
$\ro(\qrGen^{\genericRand} \bmod \rsaMod)$ only if it is able to compute it from the challenge, or if it has correctly guessed it in $\oracleQueries'$ 
queries. The latter happens with probability at most $\oracleQueries' 
/\setSize{\qrGroup}$. 

If \indsrAdversary{3} can distinguish between $\qrGen^{\left(\frac{\genericRand 		
\primeRO(\genericSetElmt[0])}{\accSet}\right)} \bmod \rsaMod$ and $\qrGen^{\left(\frac{\genericRand 
\primeRO(\genericSetElmt[1])}{\accSet}\right)} \bmod \rsaMod$ then it can also distinguish between 
$\qrGen^{\genericRand \primeRO(\genericSetElmt[0])} \bmod \rsaMod$ and $\qrGen^{\genericRand 
\primeRO(\genericSetElmt[1])} \bmod \rsaMod$ because $\indsrAdversary{3}$ knows $\accSet$.
Let $\genericRand = \modQuotient \qrGroupOrder + 
\modRemainder$. It is known that \modRemainder is uniformly distributed in 
$\integersModArg{\qrGroupOrder}$ when $\genericRand \sample \nats[1][\rsaMod^2]$~\cite{cramer2000signature}. Thus, $\genericRand 
\primeRO(\genericSetElmt[\indBit]) \bmod \qrGroupOrder$ is uniformly distributed in 
$\integersModArg{\qrGroupOrder}$ 
and $\qrGen^{\genericRand \primeRO(\genericSetElmt[0])} \bmod \rsaMod$ is statistically 
indistinguishable from $\qrGen^{\genericRand \primeRO(\genericSetElmt[1])} \bmod \rsaMod$ since
$\qrGen$ is a generator of \qrGroup. 


Finally, if $\indsrAdversary{3}$ queries the random oracle for $\qrGen^{\genericRand} \bmod 
\rsaMod$ without guessing, then it has computed $\qrGen^{\genericRand 
	\primeRO(\genericSetElmt[\indBit])^{\primeRO(\genericSetElmt[\indBit])^{-1} \bmod 
\qrGroupOrder}} \bmod \rsaMod$.  In this case, \rsaAdversary can solve the strong QR-RSA problem 
with the help of \indsrAdversaryComb. Consider the following simulation where $\rsaAdversary$ 
programs the random oracles $\ro, \primeRO'$. 

\enumListStart
%	\item Upon getting the tuple $(\rsaMod, \genericFnOutput)$ where 
%	$\genericFnOutput \in \qrGroup$, \rsaAdversary first checks that $\genericFnOutput$ is 
%	a generator of $\qrGroup$ by verifying that $\gcd(\rsaMod, \genericFnOutput - 1)  = 
%	1$~\cite{micciancio2005rsa}. Otherwise, \rsaAdversary aborts. 
%	
		\item \rsaAdversary selects \oracleQueries random prime numbers less than $\rsaMod$, 
	$\randPrimes[1], \dots, 
	\randPrimes[\oracleQueries]$ and computes $\randPrimeAccSet \gets \prod\limits_{\setIdx \in 
		[1, \oracleQueries]} \randPrimes[\setIdx]$. Then, it computes 	$\simFnOutput \gets 
	\genericFnOutput^{\randPrimeAccSet} \bmod \rsaMod$. \rsaAdversary publishes 
	$\simFnOutput$ as part of the public key. 
	
	\item\rsaAdversary selects a random index $\randIdx \in [1, \oracleQueries]$ and programs the 
	random oracle $\primeRO$ such for the $\setIdx$-th oracle query for all $\setIdx \neq \randIdx$, the random 
	oracle returns $\randPrimes[\setIdx]$. For the $\randIdx$-th query, the random oracle returns a 
	random prime number $\rsaExponent \notin \randPrimes[1], \dots, 
	\randPrimes[\oracleQueries]$
	
	%, for $\setIdx \neq \randIdx$ and returns \rsaExponent for the $\randIdx$-th  query. 
	
	\item \rsaAdversary answers the signing oracle queries from \indsrAdversary{1} and 
	\indsrAdversary{2} for the sets $\genericSetVar[1], \dots, 
	\genericSetVar[\numberOfSetsQueried]$. Specifically, upon receiving 
	$\genericSetVar$ from \indsrAdversary{1}, \rsaAdversary returns 
	$\simFnOutput^{\left(\frac{1}{\challengeSetAcc}\right)} \bmod \rsaMod$ (see the proof of \thmref{thm:unforge} to see how this is computed). If 
	$\primeRO(\genericSetElmt) = \rsaExponent$ for some $\genericSetElmt \in \genericSetVar$, then 
	the simulation aborts. 

	\item Upon receiving $\genericSetElmt[0], \genericSetElmt[1]$ from $\indsrAdversary{2}$, 
	\rsaAdversary checks that these elements are not in of any of the sets queried by  
	\indsrAdversary{1} and \indsrAdversary{2}. Note that the values of 
	$\primeRO(\genericSetElmt[0])$ and $\primeRO(\genericSetElmt[1])$ are 
	in $\randPrimes[1],\dots, \randPrimes[\randIdx-1], \rsaExponent, 
	\randPrimes[\randIdx +1], \dots, \randPrimes[\oracleQueries]$. 
	
	\item \rsaAdversary computes $\randPrimeAccSet' \gets \{ \randPrimes[i]:\randPrimes[i] \in \randPrimes[1], \dots, \randPrimes[\oracleQueries],  \forall \genericSetElmt \in \genericSetVar, \primeRO(\genericSetElmt)   \neq \randPrimes[i] \}$. 

\item \rsaAdversary computes $\dhVar[1] \gets \genericFnOutput^{\genericRandAlt P'} = 
	\simFnOutput^{ \left(\frac{\genericRandAlt}{\challengeSetAcc}\right)}$ where 
	$\genericRandAlt \sample [1, \rsaMod^2]$ and sends the challenge 
	$\tupleDefn{\dhVar[1], \genericRandAltAlt}$ to \indsrAdversary{3} where $\genericRandAltAlt \sample 
		\setDefn{0,1}^{\secParam}$. 
	
	
	\item If $\indsrAdversary{3}$ queries \ro on some value $\genericVal$ such that 
	$\genericVal^{\rsaExponent} = \simFnOutput^{\genericRandAlt} = 
	\genericFnOutput^{\genericRandAlt 
	\randPrimeAccSet'}$, then \rsaAdversary computes the values of $\euclideanConst[1], 
	\euclideanConst[2]$ using the 
	extended Euclid algorithm such that 
	$\euclideanConst[1] \rsaExponent + \euclideanConst[2] \genericRandAlt \randPrimeAccSet' = 1$. 
	Note that $\rsaExponent$ is co-prime to $\genericRandAlt \randPrimeAccSet'$ with overwhelming 
	probability. Then, we have 
	$(\genericVal^{\euclideanConst[2]} (\genericFnOutput)^{\euclideanConst[1]})^{\rsaExponent} 
	= \genericFnOutput$. \rsaAdversary outputs $\genericVal^{\euclideanConst[2]} 
	(\genericFnOutput)^{\euclideanConst[1]}$ as the $\rsaExponent$-th root of $\genericFnOutput$ 
	
\enumListEnd


To see that the simulation is indistinguishable, observe that 
$\simFnOutput^{\left(\frac{1}{\challengeSetAcc}\right)} \bmod \rsaMod$ is a valid and an 
indistinguishable signature 
(see the proof of \thmref{thm:unforge} for details). Also, as discussed before $ 	
\simFnOutput^{\left(\frac{\genericRandAlt}{\challengeSetAcc}\right)} \bmod \rsaMod$ is a uniformly 
sampled element 
in \qrGroup 
and is indistinguishable from $\simFnOutput^{\left(\frac{\genericRand 
\genericSetElmt[\indBit]}{\challengeSetAcc}\right)} \bmod \rsaMod$ because 
$\simFnOutput^{\left(\frac{
\genericSetElmt[\indBit]}{\challengeSetAcc}\right)} \bmod \rsaMod$ is a generator of \qrGroup with 
overwhelming probability. Thus, the simulation succeeds and \rsaAdversary wins when it programs 
either of the challenges $\genericSetElmt[0]$ and $\genericSetElmt[1]$ as  
$\primeRO(\genericSetElmt[0]) = \rsaExponent$ or 
$\primeRO(\genericSetElmt[1]) = \rsaExponent$, which happens with probability at least 
$1/\oracleQueries$. 


\end{proof}
\iffalse 
\subsection{Proof of robustness to key recovery}

\begin{theorem}
\label{thm:robustness}
	The signature scheme \ourSigScheme is robust against key recovery, that is 
	assuming the strong QR-RSA problem is hard to solve in \qrGroup where $\rsaMod$ is the product of safe primes.
	
	\[
	\recoverAdv \le   \frac{\oracleQueries' }{\setSize{\qrGroup}}  + 2 \Advantage{\indsrLabel}{\sigScheme}(\timeBound, 
	\oracleQueries, \oracleQueries')
	\]
Here \oracleQueries is the number of distinct queries to the random oracle \primeRO and 
$\oracleQueries'$ is the number of distinct queries to the random oracle \ro.

\end{theorem}

\begin{proof}
The proof straightforwardly reduces to the IND-SR security of the signature scheme. Since $\genericSigPair[2]$ is the output of a random oracle, \recoverAdversary{2} has either randomly guessed $\qrGen^{\genericRand}$ or computed it from $\langle \genericSetVar, \genericCInputElmt, \genericSig[\genericSetVar], \genericSig[\genericSetVar \setminus \{\genericCInputElmt\}],  \genericSigPair[1], \pubKey, \adversaryState{1} \rangle$. The former case happens with probability at most $\frac{\oracleQueries' }{\setSize{\qrGroup}} $. In the latter case, if \recoverAdversary{2} produces $\qrGen^{\genericRand}$ with probability \recoverAdv, then this adversary can be used as a subroutine by \indsrAdversary{3} to win $\Experiment{\indsrLabel[\indBit]}{\sigScheme}$ with advantage  at least $1/2 \recoverAdv$. 
\end{proof}
\fi 